% Options for packages loaded elsewhere
\PassOptionsToPackage{unicode}{hyperref}
\PassOptionsToPackage{hyphens}{url}
\PassOptionsToPackage{dvipsnames,svgnames,x11names}{xcolor}
\documentclass[
  10pt,
  fontset=fandol]{ctexart}
\usepackage{xcolor}
\usepackage[margin=16mm]{geometry}
\usepackage{amsmath,amssymb}
\setcounter{secnumdepth}{-\maxdimen} % remove section numbering
\usepackage{iftex}
\ifPDFTeX
  \usepackage[T1]{fontenc}
  \usepackage[utf8]{inputenc}
  \usepackage{textcomp} % provide euro and other symbols
\else % if luatex or xetex
  \usepackage{unicode-math} % this also loads fontspec
  \defaultfontfeatures{Scale=MatchLowercase}
  \defaultfontfeatures[\rmfamily]{Ligatures=TeX,Scale=1}
\fi
\usepackage{lmodern}
\ifPDFTeX\else
  % xetex/luatex font selection
\fi
% Use upquote if available, for straight quotes in verbatim environments
\IfFileExists{upquote.sty}{\usepackage{upquote}}{}
\IfFileExists{microtype.sty}{% use microtype if available
  \usepackage[]{microtype}
  \UseMicrotypeSet[protrusion]{basicmath} % disable protrusion for tt fonts
}{}
\makeatletter
\@ifundefined{KOMAClassName}{% if non-KOMA class
  \IfFileExists{parskip.sty}{%
    \usepackage{parskip}
  }{% else
    \setlength{\parindent}{0pt}
    \setlength{\parskip}{6pt plus 2pt minus 1pt}}
}{% if KOMA class
  \KOMAoptions{parskip=half}}
\makeatother
\usepackage{longtable,booktabs,array}
\newcounter{none} % for unnumbered tables
\usepackage{calc} % for calculating minipage widths
% Correct order of tables after \paragraph or \subparagraph
\usepackage{etoolbox}
\makeatletter
\patchcmd\longtable{\par}{\if@noskipsec\mbox{}\fi\par}{}{}
\makeatother
% Allow footnotes in longtable head/foot
\IfFileExists{footnotehyper.sty}{\usepackage{footnotehyper}}{\usepackage{footnote}}
\makesavenoteenv{longtable}
\setlength{\emergencystretch}{3em} % prevent overfull lines
\providecommand{\tightlist}{%
  \setlength{\itemsep}{0pt}\setlength{\parskip}{0pt}}
\usepackage{amsmath,amssymb,mathtools,needspace}
\xeCJKDeclareCharClass{CJK}{"2032,"0400->"04FF,"2160->"216B,"2460->"2473,"25A0,"25C7,"25CB,"25CF}
\IfFontExistsTF{Arial Unicode MS}{\xeCJKsetup{AutoFallBack=true}\setCJKfallbackfamilyfont{\CJKrmdefault}{Arial Unicode MS}}{}
\setcounter{secnumdepth}{0}
\setlength{\emergencystretch}{2em}
\allowdisplaybreaks
\usepackage{bookmark}
\IfFileExists{xurl.sty}{\usepackage{xurl}}{} % add URL line breaks if available
\urlstyle{same}
\hypersetup{
  pdftitle={幂法及反幂法},
  colorlinks=true,
  linkcolor={Maroon},
  filecolor={Maroon},
  citecolor={Blue},
  urlcolor={Blue},
  pdfcreator={LaTeX via pandoc}}

\title{幂法及反幂法}
\author{}
\date{}

\begin{document}
\maketitle

\subsection{9.2 幂法及反幂法}\label{ux5e42ux6cd5ux53caux53cdux5e42ux6cd5}

\subsubsection{9.2.1 幂法}\label{ux5e42ux6cd5}

在一些工程、物理问题中，通常只需要求出矩阵的按模最大的特征值（称为 \(A\) 的主特征值）和相应的特征向量，对于解这种特征值问题，应用幂法是合适的。

幂法是一种计算实矩阵 \(A\) 的主特征值的一种迭代法，它最大的优点是方法简单，对于稀疏矩阵较合适，但有时收敛速度很慢。

设实矩阵 \(A=(a_{ij})_n\) 有一个完全的特征向量组，其特征值为 \(\lambda_1,\lambda_2,\cdots,\lambda_n\)，相应的特征向量为 \(x_1,x_2,\cdots,x_n\)。已知 \(A\) 的主特征值是实根，且满足条件

\[
|\lambda_1|>|\lambda_2|\geq|\lambda_3|\geq\cdots\geq|\lambda_n|.
\tag{9.2.1}
\]

幂法的基本思想是任取一个非零的初始向量 \(v_0\)，由矩阵 \(A\) 构造一向量序列

\[
\left\{
\begin{aligned}
v_1&=Av_0,\\
v_2&=Av_1=A^2v_0,\\
&\quad\vdots\\
v_{k+1}&=Av_k=A^{k+1}v_0,\\
&\quad\vdots
\end{aligned}
\right.
\tag{9.2.2}
\]

称为迭代向量。由假设，\(v_0\) 可表示为

\[
v_0=a_1x_1+a_2x_2+\cdots+a_nx_n\quad\text{（设 }a_1\ne0\text{）},
\tag{9.2.3}
\]

于是

\[
\begin{aligned}
v_k&=Av_{k-1}=A^kv_0=a_1\lambda_1^kx_1+a_2\lambda_2^kx_2+\cdots+a_n\lambda_n^kx_n\\
&=\lambda_1^k\left[a_1x_1+\sum_{i=2}^na_1(\lambda_i/\lambda_1)^kx_i\right]=\lambda_1^k(a_1x_1+\varepsilon_k),
\end{aligned}
\]

其中 \(\varepsilon_k=\displaystyle\sum_{i=2}^na_1(\lambda_i/\lambda_1)^kx_i\)。由假设 \(|\lambda_i/\lambda_1|<1\ (i=2,3,\cdots,n)\)，故 \(\varepsilon_k\to0\ (k\to\infty)\)，从而

\[
\lim_{k\to\infty}\frac{v_k}{\lambda_1^k}=a_1x_1.
\tag{9.2.4}
\]

这说明序列 \(\dfrac{v_k}{\lambda_1^k}\) 越来越接近 \(A\) 的对应于 \(\lambda_1\) 的特征向量，或者说当 \(k\) 充分大时

\[
v_k\approx a_1\lambda_1^kx_1,
\tag{9.2.5}
\]

即迭代向量 \(v_k\) 为 \(\lambda_1\) 的特征向量的近似向量（除一个因子外）。

下面再考虑主特征值 \(\lambda_1\) 的计算。用 \((v_k)_i\) 表示 \(v_k\) 的第 \(i\) 个分量，则

\[
\frac{(v_{k+1})_i}{(v_k)_i}=\lambda_1\left\{\frac{a_1(x_1)_i+(\varepsilon_{k+1})_i}{a_1(x_1)_i+(\varepsilon_k)_i}\right\},
\tag{9.2.6}
\]

\begin{quote}
\textbf{校注（agent 补充，ch09a-E003）} 本页 \(v_k\) 第二行求和项及 \(\varepsilon_k\) 定义中的系数下标，原图均印为 \(a_1\)，已照录；与前一行 \(a_i\lambda_i^kx_i\) 的展开相比较，两处应为 \(a_i\) 才能保持恒等，属原书疑误。见\href{../assets/ch09a/review-p235-power-expansion.png}{放大原式}。
\end{quote}

\href{../../source-images/pdf-236.jpeg}{原页图像}

故

\[
\lim_{k\to\infty}\frac{(v_{k+1})_i}{(v_k)_i}=\lambda_1,
\tag{9.2.7}
\]

也就是说，两相邻迭代向量分量的比值收敛到主特征值。

这种由已知非零向量 \(v_0\) 及矩阵 \(A\) 的乘幂 \(A^k\) 构造向量序列 \(\{v_k\}\) 以计算 \(A\) 的主特征值 \(\lambda_1\)（利用式（9.2.7））及相应特征向量（利用式（9.2.5））的方法称为幂法。

由式（9.2.6）知，\((v_{k+1})_i/(v_k)_i\to\lambda_1\) 的收敛速度由比值 \(r=\lambda_2/\lambda_1\) 来确定，\(r\) 越小收敛越快，但当 \(r=\lambda_2/\lambda_1\approx1\) 时收敛可能就很慢。

总结上述讨论，有如下定理。

\textbf{定理 9.5} 设 \(A\in\mathbf R^{n\times n}\) 有 \(n\) 个线性无关的特征向量，主特征值 \(\lambda_1\) 满足

\[
|\lambda_1|>|\lambda_2|\geq|\lambda_3|\geq\cdots\geq|\lambda_n|,
\]

则对于任何非零初始向量 \(v_0\ (a_1\ne0)\)，式（9.2.4）、（9.2.7）成立。

设 \(A\) 的主特征值为实重根，即 \(\lambda_1=\lambda_2=\cdots=\lambda_r\)，且 \(|\lambda_r|>|\lambda_{r+1}|\geq\cdots\geq|\lambda_n|\)，又设 \(A\) 有 \(n\) 个线性无关的特征向量，\(\lambda_1\) 对应的 \(r\) 个线性无关特征向量为 \(x_1,x_2,\cdots,x_r\)，则由式（9.2.2），有

\[
v_k=A^kv_0=\lambda_1^k\left\{\sum_{i=1}^ra_ix_i+\sum_{i=r+1}^na_i\left(\frac{\lambda_i}{\lambda_1}\right)^kx_i\right\},\quad
\lim_{k\to\infty}\frac{v_k}{\lambda_1^k}=\sum_{i=1}^ra_ix_i\quad\text{（设 }\sum_{i=1}^ra_ix_i\ne0\text{）}.
\]

这说明当 \(A\) 的主特征值是实的重根时，定理 9.5 的结论还是正确的。

应用幂法计算 \(A\) 的主特征值 \(\lambda_1\) 及对应的特征向量时，如果 \(|\lambda_1|>1\)（或 \(|\lambda_1|<1\)），迭代向量 \(v_k\) 的各个不等于零的分量将随 \(k\to\infty\) 而趋于无穷（或趋于零），这样在用计算机计算时就可能``溢出''。为了克服这个缺点，就需要将迭代向量加以规范化。

设有一向量 \(v\ne0\)，将其规范化得到向量 \(u=\dfrac{v}{\max(v)}\)，其中 \(\max(v)\) 表示向量 \(v\) 的绝对值最大的分量。

在定理 9.5 的条件下幂法可这样进行：任取一初始向量 \(v_0\ne0\ (a_1\ne0)\)，构造向量序列

\[
\left\{
\begin{aligned}
v_1&=Au_0=Av_0,&u_1&=\frac{v_1}{\max(v_1)}=\frac{Av_0}{\max(Av_0)},\\
v_2&=Au_1=\frac{A^2v_0}{\max(Av_0)},&u_2&=\frac{v_2}{\max(v_2)}=\frac{A^2v_0}{\max(A^2v_0)},\\
&\quad\vdots&&\quad\vdots\\
v_k&=\frac{A^kv_0}{\max(A^{k-1}v_0)},&u_k&=\frac{A^kv_0}{\max(A^kv_0)}.
\end{aligned}
\right.
\]

由式（9.2.3），有

\[
Av_0=\sum_{i=1}^na_i\lambda_i^kx_i=\lambda_i^k\left[a_1x_1+\sum_{i=2}^na_i\left(\frac{\lambda_i}{\lambda_1}\right)^kx_i\right],
\tag{9.2.8}
\]

\begin{quote}
\textbf{校注（agent 补充，ch09a-E004）} 式（9.2.8）原图左端为 \(Av_0\)，右端方括号前为 \(\lambda_i^k\)，均照录。由 \(v_0=\sum_i a_ix_i\) 和 \(Ax_i=\lambda_ix_i\)，中间求和等于 \(A^kv_0\)，而提取的公共因子应为 \(\lambda_1^k\)；故有漏印幂次及因子下标疑误。见\href{../assets/ch09a/review-p236-9-2-8.png}{放大原式}。
\end{quote}

\href{../../source-images/pdf-237.jpeg}{原页图像}

\[
\begin{aligned}
u_k&=\frac{A^kv_0}{\max(A^kv_0)}
=\frac{\lambda_i^k\left[a_1x_1+\displaystyle\sum_{i=2}^na_i\left(\frac{\lambda_i}{\lambda_1}\right)^kx_i\right]}{\max\left[\lambda_1^k\left(a_1x_1+\displaystyle\sum_{i=2}^na_i\left(\frac{\lambda_i}{\lambda_1}\right)^kx_i\right)\right]}\\
&=\frac{\left[a_1x_1+\displaystyle\sum_{i=2}^na_i\left(\frac{\lambda_i}{\lambda_1}\right)^kx_i\right]}{\max\left[a_1x_1+\displaystyle\sum_{i=2}^na_i\left(\frac{\lambda_i}{\lambda_1}\right)^kx_i\right]}
\to\frac{x_1}{\max(x_1)}\qquad(k\to\infty).
\end{aligned}
\]

这说明规范化向量序列收敛到主特征值对应的特征向量。

同理，可得到

\[
v_k=\frac{\lambda_1^k\left[a_1x_1+\displaystyle\sum_{i=2}^na_i\left(\frac{\lambda_i}{\lambda_1}\right)^kx_i\right]}{\max\left[\lambda_1^{k-1}a_1x_1+\displaystyle\sum_{i=2}^na_i\left(\frac{\lambda_i}{\lambda_1}\right)^{k-1}x_i\right]},
\]

\[
\max(v_k)=\frac{\lambda_1\max\left[a_1x_1+\displaystyle\sum_{i=2}^na_i\left(\frac{\lambda_i}{\lambda_1}\right)^kx_i\right]}{\max\left[a_1x_1+\displaystyle\sum_{i=2}^na_i\left(\frac{\lambda_i}{\lambda_1}\right)^{k-1}x_i\right]}\to\lambda_1\qquad(k\to\infty).
\]

收敛速度由比值 \(r=\lambda_2/\lambda_1\) 确定。

总结上述讨论，有如下定理。

\textbf{定理 9.6} 设 \(A\in\mathbf R^{n\times n}\) 有 \(n\) 个线性无关的特征向量，主特征值 \(\lambda_1\) 满足 \(|\lambda_1|>|\lambda_2|\geq|\lambda_3|\geq\cdots\geq|\lambda_n|\)，则对于任意非零初始向量 \(v_0=u_0\ (a_1\ne0)\)，按下述方法构造的向量序列

\[
\left\{
\begin{aligned}
v_0&=u_0\ne0,\\
v_k&=Au_{k-1},\\
u_k&=\frac{v_k}{\max(v_k)}
\end{aligned}
\right.\qquad(k=1,2,\cdots),
\tag{9.2.9}
\]

有

\[
\lim_{k\to\infty}u_k=\frac{x_1}{\max(x_1)},\quad\lim_{k\to\infty}\max(v_k)=\lambda_1.
\]

\textbf{例 9.1} 用幂法计算 \(A=\begin{pmatrix}1.0&1.0&0.5\\1.0&1.0&0.25\\0.5&0.25&2.0\end{pmatrix}\) 的主特征值和相应的特征向量。

\textbf{解} 计算过程如表 9.1 所示。

下述结果是用 8 位浮点数字进行运算得到的，\(u_k\) 的分量值是舍入值。于是得到

\[
\lambda_1\approx2.536\,532\,3
\]

及相应特征向量 \((0.748\,2,\ 0.649\,7,\ 1)^{\mathrm T}\)。\(\lambda_1\) 和相应的特征向量真值（8 位数字）为

\[
\lambda_1=2.536\,525\,8,\quad\widetilde{x}_1=(0.748\,221\,16,\ 0.649\,661\,16,\ 1)^{\mathrm T}.
\]

\begin{quote}
\textbf{校注（agent 补充，ch09a-E005）} 页首 \(u_k\) 第一行分子方括号前原印 \(\lambda_i^k\)，分母对应因子为 \(\lambda_1^k\)，已照录；与 \(A^kv_0\) 的展开及下一行相消过程对照，分子因子应为 \(\lambda_1^k\)。见\href{../assets/ch09a/review-p237-u-k-numerator.png}{放大原式}。

\textbf{校注（agent 补充，ch09a-E006）} 本页 \(v_k\) 式分母中，原图 \(\lambda_1^{k-1}\) 仅直接乘 \(a_1x_1\)，随后为加号及求和项，已照录。由上一页 \(v_k=A^kv_0/\max(A^{k-1}v_0)\)，分母应把 \(\lambda_1^{k-1}\) 乘在整个 \(a_1x_1+\sum_{i=2}^na_i(\lambda_i/\lambda_1)^{k-1}x_i\) 上，故疑漏一层括号。见\href{../assets/ch09a/review-p237-v-k-denominator.png}{放大原式}。
\end{quote}

\href{../../source-images/pdf-238.jpeg}{原页图像}

\textbf{表 9.1}

{\def\LTcaptype{none} % do not increment counter
\begin{longtable}[]{@{}rcr@{}}
\toprule\noalign{}
\(k\) & \(u_k^{\mathrm T}\)（规范化向量） & \(\max(v_k)\) \\
\midrule\noalign{}
\endhead
\bottomrule\noalign{}
\endlastfoot
0 & \((1,1,1)\) & \\
1 & \((0.909\,1,\ 0.818\,2,\ 1)\) & \(2.750\,000\,0\) \\
5 & \((0.765\,1,\ 0.667\,4,\ 1)\) & \(2.558\,791\,8\) \\
10 & \((0.749\,4,\ 0.650\,8,\ 1)\) & \(2.538\,002\,9\) \\
15 & \((0.748\,3,\ 0.649\,7,\ 1)\) & \(2.536\,625\,6\) \\
16 & \((0.748\,3,\ 0.649\,7,\ 1)\) & \(2.536\,584\,0\) \\
17 & \((0.748\,2,\ 0.649\,7,\ 1)\) & \(2.536\,559\,8\) \\
18 & \((0.748\,2,\ 0.649\,7,\ 1)\) & \(2.536\,545\,6\) \\
19 & \((0.748\,2,\ 0.649\,7,\ 1)\) & \(2.536\,537\,4\) \\
20 & \((0.748\,2,\ 0.649\,7,\ 1)\) & \(2.536\,532\,3\) \\
\end{longtable}
}

\subsubsection{9.2.2 加速方法}\label{ux52a0ux901fux65b9ux6cd5}

\paragraph{1. 原点平移法}\label{ux539fux70b9ux5e73ux79fbux6cd5}

由前面讨论知道，应用幂法计算 \(A\) 的主特征值时，其收敛速度主要由比值 \(r=\dfrac{\lambda_1}{\lambda_2}\) 来决定，但当 \(r\) 接近于 1 时，收敛可能很慢。这时，一个补救的办法是采用加速收敛的方法。

引进矩阵 \(B=A-pI\)，其中 \(p\) 为选择参数。

设 \(A\) 的特征值为 \(\lambda_1,\lambda_2,\cdots,\lambda_n\)，则 \(B\) 的相应特征值为 \(\lambda_1-p,\lambda_2-p,\cdots,\lambda_n-p\)，而且 \(A,B\) 的特征向量相同。

如果需要计算 \(A\) 的主特征值 \(\lambda_1\)，就要选择适当的 \(p\) 使 \(\lambda_1-p\) 仍然是 \(B\) 的主特征值，且使

\[
\left|\frac{\lambda_2-p}{\lambda_1-p}\right|<\left|\frac{\lambda_2}{\lambda_1}\right|.
\]

对 \(B\) 应用幂法，使得在计算 \(B\) 的主特征值 \(\lambda_1-p\) 的过程中得到加速。这种方法通常称为原点平移法。对于 \(A\) 的特征值的某种分布，它是十分有效的。

\textbf{例 9.2} 设 \(A=(a_{ij})_4\) 有特征值 \(\lambda_j=15-j\quad(j=1,2,3,4)\)，比值 \(r=\lambda_2/\lambda_1\approx0.9\)。

作变换

\[
B=A-pI\quad(p=12),
\]

则 \(B\) 的特征值为 \(\mu_1=2,\ \mu_2=1,\ \mu_3=0,\ \mu_4=-1\)。应用幂法计算 \(B\) 的主特征值 \(\mu_1\) 的收敛速度的比值为

\[
\left|\frac{\mu_2}{\mu_1}\right|=\left|\frac{\lambda_2-p}{\lambda_1-p}\right|=\frac12<\left|\frac{\lambda_2}{\lambda_1}\right|\approx0.9.
\]

\begin{quote}
\textbf{校注（agent 补充，ch09a-E007）} 本页``原点平移法''首段原印 \(r=\lambda_1/\lambda_2\)，已照录；前页及本页例 9.2 均用 \(r=\lambda_2/\lambda_1\)，与幂法误差项 \((\lambda_i/\lambda_1)^k\) 对照，此处疑将分子、分母颠倒。见\href{../assets/ch09a/review-p238-ratio.png}{放大原式}。
\end{quote}

\href{../../source-images/pdf-239.jpeg}{原页图像}

虽然常常能够选择有利的 \(p\) 值，使幂法得到加速，但设计一个自动选择适当参数 \(p\) 的过程是困难的。

下面考虑当 \(A\) 的特征值是实数时，怎样选择 \(p\) 使用幂法计算 \(\lambda_1\) 以得到加速。

设 \(A\) 的特征值满足

\[
\lambda_1>\lambda_2\geq\cdots\geq\lambda_{n-1}>\lambda_n,
\tag{9.2.10}
\]

则不管 \(p\) 如何选择，\(B=A-pI\) 的主特征值为 \(\lambda_1-p\) 或 \(\lambda_n-p\)。当希望计算 \(\lambda_1\) 及 \(x_1\) 时，首先应选择 \(p\) 使 \(|\lambda_1-p|>|\lambda_n-p|\)，且使收敛速度的比值

\[
\omega=\max\left\{\frac{|\lambda_2-p|}{|\lambda_1-p|},\frac{|\lambda_n-p|}{|\lambda_1-p|}\right\}=\min,
\]

显然，当 \(\lambda_2-p=-(\lambda_n-p)\)，\(p=\dfrac{\lambda_2+\lambda_n}{2}\equiv p^*\) 时 \(\omega\) 为最小，这时收敛速度的比值为

\[
\frac{\lambda_2-p^*}{\lambda_1-p^*}=-\frac{\lambda_n-p^*}{\lambda_1-p^*}\equiv\frac{\lambda_2-\lambda_n}{2\lambda_1-\lambda_2-\lambda_n}.
\]

当 \(A\) 的特征值满足式（9.2.10）且 \(\lambda_2,\lambda_n\) 能初步估计时，就能确定 \(p^*\) 的近似值。

当希望计算 \(\lambda_n\) 时，应选择 \(p=\dfrac{\lambda_1+\lambda_{n-1}}{2}=p^*\)，使得应用幂法计算 \(\lambda_n\) 得到加速。

\textbf{例 9.3} 计算例 9.1 矩阵 \(A\) 的主特征值。

\textbf{解} 作变换 \(B=A-pI\)，取 \(p=0.75\)，则

\[
B=\begin{pmatrix}
0.25&1&0.5\\
1&0.25&0.25\\
0.5&0.25&1.25
\end{pmatrix}.
\]

对 \(B\) 应用幂法，计算结果如表 9.2 所示。

\textbf{表 9.2}

{\def\LTcaptype{none} % do not increment counter
\begin{longtable}[]{@{}rcr@{}}
\toprule\noalign{}
\(k\) & \(u_k^{\mathrm T}\)（规范化向量） & \(\max(v_k)\) \\
\midrule\noalign{}
\endhead
\bottomrule\noalign{}
\endlastfoot
0 & \((1,1,1)\) & \\
5 & \((0.751\,6,\ 0.652\,2,\ 1)\) & \(1.791\,401\,1\) \\
6 & \((0.749\,1,\ 0.651\,1,\ 1)\) & \(1.788\,844\,3\) \\
7 & \((0.748\,8,\ 0.650\,1,\ 1)\) & \(1.787\,330\,0\) \\
8 & \((0.748\,4,\ 0.649\,9,\ 1)\) & \(1.786\,915\,2\) \\
9 & \((0.748\,3,\ 0.649\,7,\ 1)\) & \(1.786\,658\,7\) \\
10 & \((0.748\,2,\ 0.649\,7,\ 1)\) & \(1.786\,591\,4\) \\
\end{longtable}
}

由此得 \(B\) 的主特征值为 \(\mu_1\approx1.786\,591\,4\)，\(A\) 的主特征值 \(\lambda_1\) 为

\[
\lambda_1\approx\mu_1+0.75=2.536\,591\,4.
\]

这个结果比例 9.1 迭代 15 次得到的结果还要好。若迭代 15 次，\(\mu_1=1.786\,525\,8\)（相应的 \(\lambda_1=2.536\,525\,8\)）。

\href{../../source-images/pdf-240.jpeg}{原页图像}

原点位移的加速方法，是一个矩阵变换方法。这种变换容易计算，又不破坏矩阵 \(A\) 的稀疏性，但 \(p\) 的选择依赖于对 \(A\) 的特征值分布的大致了解。

\paragraph{2. Rayleigh 商加速法}\label{rayleigh-ux5546ux52a0ux901fux6cd5}

由定理 9.4 知，对称矩阵 \(A\) 的 \(\lambda_1\) 及 \(\lambda_n\) 可用 Rayleigh 商的极值来表示。下面将把 Rayleigh 商应用到用幂法计算实对称矩阵 \(A\) 的主特征值的加速收敛上来。

\textbf{定理 9.7} 设 \(A\in\mathbf R^{n\times n}\) 为对称矩阵，特征值满足 \(|\lambda_1|>|\lambda_2|\geq|\lambda_3|\geq\cdots\geq|\lambda_n|\)，对应的特征向量满足 \((x_i,x_j)=\delta_{ij}\)，应用幂法（式（9.2.9））计算 \(A\) 的主特征值 \(\lambda_1\)，则规范化向量 \(u_k\) 的 Rayleigh 商给出 \(\lambda_1\) 的较好的近似，即

\[
\frac{(Au_k,u_k)}{(u_k,u_k)}=\lambda_1+O\left(\left(\frac{\lambda_2}{\lambda_1}\right)^{2k}\right).
\]

\textbf{证明} 由式（9.2.8）及 \(u_k=\dfrac{A^ku_0}{\max(A^ku_0)}\)，\(v_{k+1}=Au_k=\dfrac{A^{k+1}u_0}{\max(A^ku_0)}\)，得

\[
\frac{(Au_k,u_k)}{(u_k,u_k)}=\frac{(A^{k+1}u_0,A^ku_0)}{(A^ku_0,A^ku_0)}=\frac{\displaystyle\sum_{j=1}^na_j^2\lambda_j^{2k+1}}{\displaystyle\sum_{j=1}^na_j^2\lambda_j^{2k}}=\lambda_1+O\left(\left(\frac{\lambda_2}{\lambda_1}\right)^{2k}\right).
\tag{9.2.11}
\]

\subsubsection{9.2.3 反幂法}\label{ux53cdux5e42ux6cd5}

反幂法用来计算矩阵按模最小的特征值及其特征向量，及计算对应于一个给定近似特征值的特征向量。

设 \(A\in\mathbf R^{n\times n}\) 为非奇异矩阵，\(A\) 的特征值依次记作 \(|\lambda_1|\geq|\lambda_2|\geq\cdots\geq|\lambda_n|\)，相应的特征向量为 \(x_1,x_2,\cdots,x_n\)，则 \(A^{-1}\) 的特征值为 \(\left|\dfrac1{\lambda_n}\right|\geq\left|\dfrac1{\lambda_{n-1}}\right|\geq\cdots\geq\left|\dfrac1{\lambda_1}\right|\)，对应的特征向量为 \(x_n,x_{n-1},\cdots,x_1\)。

因此，计算 \(A\) 的按模最小的特征值 \(\lambda_n\) 的问题就是计算 \(A^{-1}\) 的按模最大的特征值问题。

对 \(A^{-1}\) 应用幂法迭代法（称为反幂法），可求得矩阵 \(A^{-1}\) 的主特征值 \(1/\lambda_n\)，从而求得 \(A\) 的按模最小的特征值 \(\lambda_n\)。

反幂法迭代公式为任取初始向量 \(v_0=u_0\ne0\)，构造向量序列

\[
\left\{
\begin{aligned}
v_k&=A^{-1}u_{k-1},\\
u_k&=\frac{v_k}{\max(v_k)}
\end{aligned}
\right.\qquad(k=1,2,\cdots).
\]

迭代向量 \(v_k\) 可以通过解方程组 \(Av_k=u_{k-1}\) 求得。

\textbf{定理 9.8} 设

\(1^\circ\) \(A\) 有 \(n\) 个线性无关的特征向量，

\(2^\circ\) \(A\) 为非奇异矩阵且其特征值满足

\href{../../source-images/pdf-241.jpeg}{原页图像}

\[
|\lambda_1|\geq|\lambda_2|\geq\cdots\geq|\lambda_{n-1}|>|\lambda_n|>0,
\]

则对任何初始非零向量 \(u_0=v_0\ (a_n\ne0)\)，由反幂法构造的向量序列 \(\{v_k\},\{u_k\}\) 满足

\[
\lim_{k\to\infty}u_k=\frac{x_n}{\max(x_n)},\quad\lim_{k\to\infty}\max(v_k)=\frac1{\lambda_n}.
\]

收敛速度的比值为 \(\left|\dfrac{\lambda_n}{\lambda_{n-1}}\right|\)。

在反幂法中也可以用原点平移法来加速迭代过程或求其他特征值及特征向量。

如果矩阵 \((A-pI)^{-1}\) 存在，显然其特征值为 \(\dfrac1{\lambda_1-p},\dfrac1{\lambda_2-p},\cdots,\dfrac1{\lambda_n-p}\)，对应的特征向量仍然是 \(x_1,x_2,\cdots,x_n\)。现对矩阵 \((A-pI)^{-1}\) 应用幂法，得到反幂法的迭代公式

\[
\left\{
\begin{aligned}
u_0&=v_0\ne0\quad\text{（初始向量）},\\
v_k&=(A-pI)^{-1}u_{k-1},\\
u_k&=\frac{v_k}{\max(v_k)}
\end{aligned}
\right.\qquad(k=1,2,\cdots).
\tag{9.2.12}
\]

如果 \(p\) 是 \(A\) 的特征值 \(\lambda_j\) 的一个近似值，且 \(|\lambda_j-p|<|\lambda_i-p|\ (i\ne j)\)，就是说，\(\dfrac1{\lambda_j-p}\) 是 \((A-pI)^{-1}\) 的主特征值，可用反幂法式（9.2.12）计算其特征值及特征向量。

设 \(A\in\mathbf R^{n\times n}\) 有 \(n\) 个线性无关的特征向量 \(x_1,x_2,\cdots,x_n\)，则

\[
u_0=\sum_{i=1}^na_ix_i\quad(a_i\ne0),\quad
v_k=\frac{(A-pI)^{-k}u_0}{\max((A-pI)^{-(k-1)}u_0)},
\]

\[
u_k=\frac{(A-pI)^{-k}u_0}{\max((A-pI)^{-k}u_0)},
\]

其中

\[
(A-pI)^{-k}u_0=\sum_{i=1}^na_i(\lambda_i-p)^{-k}x_i.
\]

\textbf{定理 9.9} 设

\(1^\circ\) \(A\in\mathbf R^{n\times n}\) 有 \(n\) 个线性无关的特征向量，\(A\) 的特征值及对应的特征向量记作 \(\lambda_i\) 及 \(x_i\ (i=1,2,\cdots,n)\)；

\(2^\circ\) \(p\) 为 \(\lambda_j\) 的近似值，\((A-pI)^{-1}\) 存在，且 \(|\lambda_j-p|<|\lambda_i-p|\ (i\ne j)\)；

\(3^\circ\) \(u_0=\displaystyle\sum_{i=1}^na_ix_i\ne0\) 为给定的初始向量 \((a_i\ne0)\)，

则由反幂法迭代公式（9.2.12）构造的向量序列 \(\{v_k\},\{u_k\}\) 满足

\[
\lim_{k\to\infty}u_k=\frac{x_j}{\max(x_j)},
\]

\[
\lim_{k\to\infty}\max(v_k)=\frac1{\lambda_j-p},\quad\text{即}\quad p+\frac1{\max(v_k)}\to\lambda_j\quad\text{（当 }k\to\infty\text{）},
\]

且收敛速度由比值 \(r=\displaystyle\max_{i\ne j}\left|\dfrac{\lambda_j-p}{\lambda_i-p}\right|\) 确定。

\href{../../source-images/pdf-242.jpeg}{原页图像}

由定理 9.9 知，对 \(A-pI\)（其中 \(p\approx\lambda_j\)）应用反幂法，可计算特征向量 \(x_j\)。只要选择的 \(p\) 是 \(\lambda_j\) 的一个较好的近似且特征值分离情况较好，一般 \(r\) 很小，常常只要迭代一两次就可完成特征向量的计算。

反幂法迭代公式中的 \(v_k\) 是通过解方程组 \((A-pI)v_k=u_{k-1}\) 求得的。为了节省工作量，可以先将 \((A-pI)\) 进行三角分解，即

\[
P(A-pI)=LU,
\]

其中 \(P\) 为某个置换矩阵，于是求 \(v_k\) 相当于解两个三角形方程组 \(Ly_k=Pu_{k-1},Uv_k=y_k\)。

反幂法迭代公式可写为

\[
\left\{
\begin{aligned}
Ly_k&=Pu_{k-1},\\
Uv_k&=y_k,\\
u_k&=\frac{v_k}{\max(v_k)}
\end{aligned}
\right.\qquad(k=1,2,\cdots).
\tag{9.2.13}
\]

实验表明，按下述方法选择 \(v_0=u_0\) 是较好的：选 \(u_0\) 使

\[
Uv_1=L^{-1}Pu_0=(1,1,\cdots,1)^{\mathrm T},
\tag{9.2.14}
\]

用回代求解式（9.2.14）即得 \(v_1\)，然后再按式（9.2.13）迭代。

\textbf{例 9.4} 用反幂法求

\[
A=\begin{pmatrix}
2&1&0\\
1&3&1\\
0&1&4
\end{pmatrix}
\]

的对应计算特征值 \(\lambda=1.267\,9\)（精确特征值为 \(\lambda_3=3-\sqrt3\)）的特征向量（用 5 位浮点数进行运算）。

\textbf{解} 用部分选主元的三角分解将 \(A-pI\)（其中 \(p=1.267\,9\)）分解为

\[
P(A-pI)=LU,
\]

其中

\[
P=\begin{pmatrix}
0&1&0\\
0&0&1\\
1&0&0
\end{pmatrix},
\]

\[
L=\begin{pmatrix}
1&0&0\\
0&1&0\\
0.732\,1&-0.268\,07&1
\end{pmatrix},\quad
U=\begin{pmatrix}
1&1.732\,1&1\\
0&1&2.732\,1\\
0&0&0.294\,05\times10^{-3}
\end{pmatrix}.
\]

由 \(Uv_1=(1,1,1)^{\mathrm T}\)，得

\[
v_1=(12\,692,\ -9\,290.3,\ 3\,400.8)^{\mathrm T},\quad
u_1=(1,\ -0.731\,98,\ 0.267\,95)^{\mathrm T},
\]

由 \(LUv_2=Pu_1\)，得

\[
v_2=(20\,404,\ -14\,937,\ 5\,467.4)^{\mathrm T},\quad
u_2=(1,\ -0.732\,06,\ 0.267\,96)^{\mathrm T}.
\]

\(\lambda_3\) 对应的特征向量是

\[
x_3=(1,\ 1-\sqrt3,\ 2-\sqrt3)^{\mathrm T}\approx(1,\ -0.732\,05,\ 0.267\,95)^{\mathrm T}.
\]

\href{../../source-images/pdf-243.jpeg}{原页图像}

由此可以看出，\(u_2\) 是 \(x_3\) 的相当好的近似。

\subsection{本节覆盖原页的校注记录}\label{ux672cux8282ux8986ux76d6ux539fux9875ux7684ux6821ux6ce8ux8bb0ux5f55}

下列记录按原页关联，含同页相邻小节的校注；计算时同时读取上文校注和完整原页。

\begin{itemize}
\tightlist
\item
  \texttt{ch09a-E003}：原书 PDF 页 235
\item
  \texttt{ch09a-E004}：原书 PDF 页 236
\item
  \texttt{ch09a-E005}：原书 PDF 页 237
\item
  \texttt{ch09a-E006}：原书 PDF 页 237
\item
  \texttt{ch09a-E007}：原书 PDF 页 238
\end{itemize}

\end{document}
